\documentclass[10pt,a4paper]{amsart}
\usepackage[margin=19mm]{geometry}
\usepackage{amsmath,amssymb,mathtools}
\usepackage[scr=rsfs,cal=euler]{mathalfa}
\usepackage{xfrac}
\usepackage[unicode,hidelinks,bookmarks=false]{hyperref}
\newcommand{\ie}{i.e.\ }
\newcommand{\cf}{cf.\ }
\newcommand{\resp}{respectively\ }
\newcommand{\Subsection}{\subsection}
\providecommand{\nobreakdash}{\nobreak\hbox{-}}
\setlength{\emergencystretch}{2em}
\multlinegap0pt
\usepackage[shortlabels]{enumitem}
\usepackage{amssymb}
\usepackage{tikz}
\newtheorem{theorem}{Theorem}
\newcommand{\N}{\mathbb{N}}   % set of natural numbers
\newcommand{\Z}{\mathbb{Z}}   % set of integer numbers
\title{On notions of expansivity for operators\\ on locally convex spaces}
\author{Nilson C. Bernardes Jr., F\'elix Mart\'inez-Gim\'enez, Francisco Rodenas}
\date{Source: arXiv:2509.12410v1 (CC BY 4.0)}
\begin{document}
\maketitle

\noindent\emph{Source and licence.} On notions of expansivity for operators\\ on locally convex spaces. Version arXiv:2509.12410v1 (CC BY 4.0), distributed by arXiv under the Creative Commons Attribution 4.0 International licence, http://creativecommons.org/licenses/by/4.0/, as recorded in the arXiv licence metadata for this exact version and shown on the abstract page https://arxiv.org/abs/2509.12410v1. Full licence notice: \texttt{licenses/arxiv-2509.12410-CC-BY-4.0.txt}.\par\smallskip

\noindent\textit{Editorial scope.} Counted unit: the article's theorem UEWS-T1 (source0.tex lines 784--814). The article's complete original proof of it is delivered with it (source0.tex lines 816--944, one \texttt{proof} environment of 129 lines). No mathematics is written, repaired or re-derived: the statement and the proof are the article's own lines, in the article's own notation, and no retained line is substituted or re-flowed.  Retained uncounted as the source-local context the proof needs: lines 520--522; lines 525--527.  Nothing was omitted from any retained span, so the package carries no omission record.  No retained line contains a citation, so the wrapper carries no bibliography and no bibliography entry.  No retained line contains a figure environment, an image-inclusion command or drawing code, so no image is delivered and the package declares no asset dependency.  Editorial formatting only, in addition: an AMS \texttt{amsart} preamble that restates the article's own declarations and macros used by the retained lines, namely the environments \texttt{theorem} on separate counters, together with the standard packages those lines need.  The retained environments print in this document as Theorem 1 in order of appearance, whereas the article prints the same items with the numbering of its own full document.  The complete native source of the article as distributed in the arXiv e-print for this exact version is bundled unchanged as \texttt{sources/185119/source0.tex}.
\begin{equation}\label{Fw-Defined}
\sup_{j \in \Z} \frac{a_{j+1,k} |w_j|}{a_{j,\ell}} < \infty.
\end{equation}

\begin{equation}\label{Fw-Invertible}
\sup_{j \in \Z} \frac{a_{j,k}}{a_{j+1,\ell} |w_j|} < \infty.
\end{equation}

\begin{theorem}\label{UEWS-T1}
Consider a K\"othe sequence space $X = \lambda_p(A,\Z)$, where $A = (a_{j,k})$ is a K\"othe matrix on $\Z$ 
and $p \in \{0\} \cup [1,\infty)$. For each $k \in \N$, define
\[
I_k = \{j \in \Z : a_{j,k} \neq 0\}.
\]
Let $w = (w_j)_{j \in \Z}$ be a sequence of nonzero scalars such that the bilateral weighted forward shift
\[
F_w : (x_j)_{j \in \Z} \in X \mapsto (w_{j-1}x_{j-1})_{j \in \Z} \in X
\]
is an invertible operator on $X$ (i.e., conditions (\ref{Fw-Defined}) and (\ref{Fw-Invertible}) hold). 
Then $F_w$ is uniformly expansive if and only if one of the following properties holds:
\begin{itemize}
\item [\rm (A)] For each $k \in \N$, there exists $\ell \in \N$ such that
    \[
    \lim_{n \to \infty} \Big(\inf_{j \in I_k} \frac{a_{j+n,\ell} |w_j \cdots w_{j+n-1}|}{a_{j,k}}\Big) = \infty.
    \]
\item [\rm (B)] For each $k \in \N$, there exists $\ell \in \N$ such that
    \[
    \lim_{n \to \infty} \Big(\inf_{j \in I_k} \frac{a_{j-n,\ell}}{a_{j,k} |w_{j-n} \cdots w_{j-1}|}\Big) = \infty.
    \]
\item [\rm (C)] For each $k \in \N$, there exists $\ell \in \N$ such that
    \[
    \lim_{n \to \infty} \Big(\inf_{j \in I_k \cap \N} \frac{a_{j+n,\ell} |w_j \cdots w_{j+n-1}|}{a_{j,k}}\Big) = \infty
    \]
    and
    \[
    \lim_{n \to \infty} \Big(\inf_{j \in I_k \cap (-\N)} \frac{a_{j-n,\ell}}{a_{j,k} |w_{j-n} \cdots w_{j-1}|}\Big) = \infty.
    \]
\end{itemize}
\end{theorem}

\begin{proof}
($\Rightarrow$): By the definition of uniform expansivity, for each $k \in \N$, there exist $\ell_k \in \N$ 
and subsets $S_k^+$ and $S_k^-$ of $S_{\|\cdot\|_k}$ such that $S_{\|\cdot\|_k} = S_k^+ \cup S_k^-$,
\[
\|F_w^n(x)\|_{\ell_k} \to \infty \text{ uniformly on } S_k^+ \text{ as } n \to \infty
\]
and
\[
\|F_w^{-n}(x)\|_{\ell_k} \to \infty \text{ uniformly on } S_k^- \text{ as } n \to \infty.
\]
For each $n \in \N$, let
\[
C_{k,n} = \inf_{x \in S_k^+} \|F_w^n(x)\|_{\ell_k} \ \ \text{ and } \ \ D_{k,n} = \inf_{x \in S_k^-} \|F_w^{-n}(x)\|_{\ell_k}.
\]
Then $C_{k,n} \to \infty$ and $D_{k,n} \to \infty$ as $n \to \infty$. Consider the sets
\[
I_k^+ = \Big\{j \in I_k : \frac{e_j}{a_{j,k}} \in S_k^+\Big\} \ \ \text{ and } \ \ I_k^- = \Big\{j \in I_k : \frac{e_j}{a_{j,k}} \in S_k^-\Big\}.
\]
Clearly, $I_k = I_k^+ \cup I_k^-$. It follows from the definitions that, for each $n \in \N$,
\[
\inf_{j \in I_k^+} \frac{a_{j+n,\ell_k} |w_j \cdots w_{j+n-1}|}{a_{j,k}} 
  = \inf_{j \in I_k^+} \Big\|F_w^n\Big(\frac{e_j}{a_{j,k}}\Big)\Big\|_{\ell_k} \geq C_{k,n}
\]
and
\[
\inf_{j \in I_k^-} \frac{a_{j-n,\ell_k}}{a_{j,k} |w_{j-n} \cdots w_{j-1}|} 
  = \inf_{j \in I_k^-} \Big\|F_w^{-n}\Big(\frac{e_j}{a_{j,k}}\Big)\Big\|_{\ell_k} \geq D_{k,n}.
\]
Therefore, 
\begin{equation}\label{UEWS-F1}
\lim_{n \to \infty} \Big(\inf_{j \in I_k^+} \frac{a_{j+n,\ell_k} |w_j \cdots w_{j+n-1}|}{a_{j,k}}\Big) = \infty\, \text{ and }
\lim_{n \to \infty} \Big(\inf_{j \in I_k^-} \frac{a_{j-n,\ell_k}}{a_{j,k} |w_{j-n} \cdots w_{j-1}|}\Big) = \infty.
\end{equation}
If $I_k^- = \emptyset$ for all $k \in \N$, then (A) holds, and if $I_k^+ = \emptyset$ for all $k \in \N$, then (B) holds.
Suppose that there exist $r,s \in \N$ such that $I_r^- \neq \emptyset$ and $I_s^+ \neq \emptyset$.
Then (\ref{UEWS-F1}) implies that there is $t \in \N$ such that $I_k = \Z$ for all $k \geq t$.
Let $k_0 = \max\{t,\ell_r,\ell_s\}$. We claim that
\begin{equation}\label{UEWS-F2}
I_k^+ \text{ is bounded below and } I_k^- \text{ is bounded above for all } k \geq k_0.
\end{equation}
Indeed, take $j_1 \in I_r^-$ and $j_2 \in I_s^+$. Then,
\[
\lim_{n \to \infty} \frac{a_{j_1-n,\ell_r}}{a_{j_1,r} |w_{j_1-n} \cdots w_{j_1-1}|} = \infty \ \ \text{ and } \
\lim_{n \to \infty} \frac{a_{j_2+n,\ell_s} |w_{j_2} \cdots w_{j_2+n-1}|}{a_{j_2,s}} = \infty,
\]
which implies that
\begin{equation}\label{UEWS-F3}
\lim_{n \to \infty} \frac{a_{-n,\ell_r}}{|w_{-n} \cdots w_{-1}|} = \infty \ \ \text{ and } \
\lim_{n \to \infty} a_{n,\ell_s} |w_0 \cdots w_{n-1}| = \infty.
\end{equation}
Assume $k \geq k_0$. If $n \in \N$ and $-n \in I_k^+$, then
\[
\frac{a_{0,\ell_k} |w_{-n} \cdots w_{-1}|}{a_{-n,\ell_r}}   \geq \frac{a_{0,\ell_k} |w_{-n} \cdots w_{-1}|}{a_{-n,k}} \geq C_{k,n}.
\]
Hence, the first equality in (\ref{UEWS-F3}) implies that $I_k^+$ must be bounded below.
Similarly, if $n \in \N$ and $n \in I_k^-$, then
\[
\frac{a_{0,\ell_k}}{a_{n,\ell_s} |w_0 \cdots w_{n-1}|} \geq \frac{a_{0,\ell_k}}{a_{n,k} |w_0 \cdots w_{n-1}|} \geq D_{k,n}.
\]
Thus, the second equality in (\ref{UEWS-F3}) implies that $I_k^-$ must be bounded above, completing the proof of our claim.

For $k \geq k_0$, (\ref{UEWS-F2}) implies that the sets $I_k^+$ and $I_k^-$ are nonempty (because $I_k = \Z$) and, consequently, 
the limits in (\ref{UEWS-F1}) do not change if we remove (resp.\ add) a finite number of integers from (resp.\ to) these sets.
Thus, it follows from (\ref{UEWS-F2}) that (\ref{UEWS-F1}) is equivalent to the limits in (C).

For $k < k_0$, we have that $I_k \subset I_{k_0}$ and, therefore, the limits in (C) hold with $\ell = \ell_{k_0}$.

\smallskip\noindent
($\Leftarrow$): We first assume that (C) holds. Given $k \in \N$, let $\ell \in \N$ be as in (C). Let 
\[
I_k^+ = I_k \cap \N_0 \ \ \text{ and } \ \ I_k^- = I_k \cap (-\N).
\]
For each $x = (x_j)_{j \in \Z} \in X$, we define $x^+ = (x^+_j)_{j \in \Z} \in X$ and $x^- = (x^-_j)_{j \in \Z} \in X$ by
\[
x^+_j = \left\{\begin{array}{cl} x_j & \text{if } j \in I_k^+\\ 0 & \text{otherwise} \end{array}\right. \ \ \text{ and } \ \ \
x^-_j = \left\{\begin{array}{cl} x_j & \text{if } j \in I_k^-\\ 0 & \text{otherwise} \end{array}\right..
\]
Let 
\[
S_k^+ = \{x \in S_{\|\cdot\|_k} : \|x^+\|_k \geq 1/2\} \ \ \text{ and } \ \ S_k^- = \{x \in S_{\|\cdot\|_k} : \|x^-\|_k \geq 1/2\}.
\]
Clearly, $S_{\|\cdot\|_k} = S_k^+ \cup S_k^-$. For each $n \in \N$, let
\[
C_n = \inf_{j \in I_k^+} \frac{a_{j+n,\ell} |w_j \cdots w_{j+n-1}|}{a_{j,k}} \ \ \text{ and } \ \
D_n = \inf_{j \in I_k^-} \frac{a_{j-n,\ell}}{a_{j,k} |w_{j-n} \cdots w_{j-1}|}\cdot
\]
By our choice of $\ell$, $C_n \to \infty$ and $D_n \to \infty$ as $n \to \infty$. Take $y \in S_k^+$. If $p \in [1,\infty)$, then
\[
\|F_w^n(y)\|_\ell^p \geq \|F_w^n(y^+)\|_\ell^p = \sum_{j \in I_k^+ + n} |a_{j,\ell}\, w_{j-n} \cdots w_{j-1}\, y^+_{j-n}|^p
  \geq C_n^p \, \|y^+\|_k^p \geq \Big(\frac{C_n}{2}\Big)^p.
\]
If $p = 0$, then
\[
\|F_w^n(y)\|_\ell \geq \|F_w^n(y^+)\|_\ell = \sup_{j \in I_k^+ + n} |a_{j,\ell}\, w_{j-n} \cdots w_{j-1}\, y^+_{j-n}|
  \geq C_n \, \|y^+\|_k \geq \frac{C_n}{2}\cdot
\]
In any case, it follows that 
\[
\|F_w^n(x)\|_\ell \to \infty \ \text{ uniformly on } S_k^+ \text{ as } n \to \infty.
\]
Now, take $z \in S_k^-$. If $p \in [1,\infty)$, then
\[
\|F_w^{-n}(z)\|_\ell^p \geq \|F_w^{-n}(z^-)\|_\ell^p 
  = \sum_{j \in I_k^- - n} \Big|\frac{a_{j,\ell}\, z^-_{j+n}}{w_j \cdots w_{j+n-1}}\Big|^p
  \geq D_n^p \, \|z^-\|_k^p \geq \Big(\frac{D_n}{2}\Big)^p.
\]
If $p = 0$, then
\[
\|F_w^{-n}(z)\|_\ell \geq \|F_w^{-n}(z^-)\|_\ell 
  = \sup_{j \in I_k^- - n} \Big|\frac{a_{j,\ell}\, z^-_{j+n}}{w_j \cdots w_{j+n-1}}\Big|
  \geq D_n \, \|z^-\|_k \geq \frac{D_n}{2}\cdot
\]
In any case, we see that 
\[
\|F_w^{-n}(x)\|_\ell \to \infty \ \text{ uniformly on } S_k^- \text{ as } n \to \infty.
\]

If (A) holds, then similar computations show that
\[
\|F_w^n(x)\|_\ell \to \infty \ \text{ uniformly on } S_{\|\cdot\|_k} \text{ as } n \to \infty.
\]

Analogously, if (B) holds, then
\[
\|F_w^{-n}(x)\|_\ell \to \infty \ \text{ uniformly on } S_{\|\cdot\|_k} \text{ as } n \to \infty.
\]

This completes the proof that $F_w$ is uniformly expansive.
\end{proof}

\end{document}
